\documentclass[11pt]{article}

\usepackage[margin=1in]{geometry}
\usepackage[T1]{fontenc}
\usepackage[utf8]{inputenc}
\usepackage{amsmath,amssymb,amsthm,mathtools}
\usepackage{enumitem}
\usepackage{hyperref}
\usepackage{longtable}

\hypersetup{
  colorlinks=true,
  linkcolor=blue,
  urlcolor=blue
}

\newtheorem{definition}{Definition}[section]
\newtheorem{theorem}[definition]{Theorem}
\newtheorem{lemma}[definition]{Lemma}
\newtheorem{proposition}[definition]{Proposition}
\newtheorem{conjecture}[definition]{Conjecture}
\newtheorem{warning}[definition]{Warning}

\DeclareMathOperator{\Hom}{Hom}
\DeclareMathOperator{\End}{End}
\DeclareMathOperator{\Irr}{Irr}

\title{Formal Definitions, Results, and Provenance for Categorical Fine Graining}
\author{CFT Anyons Formalisation Workspace}
\date{May 14, 2026}

\begin{document}
\maketitle

\begin{abstract}
This document is the self-contained mathematical record for the
formalisation project. Every external mathematical input is tied to a local
source file. Every result is labelled as proved, ready for adversarial
verification, in progress, or conjectural. A result is not accepted merely
because it appears in the handoff document; acceptance requires local source
provenance, a proof-tree entry in the Adversarial Proof Framework, and where
applicable three independent checks in Lean, Julia, and WolframScript.
\end{abstract}

\tableofcontents

\section{Status and Ground Rules}

\begin{definition}[Acceptance gate]
A statement is accepted only after all applicable items below have been
recorded.
\begin{enumerate}[label=\arabic*.]
\item A local source string locator for every external definition, equation,
or claim.
\item A proof-tree node in the Adversarial Proof Framework.
\item A Lean proof whenever the claim is within the current formal scope.
\item A Julia check for numerical or finite examples.
\item A WolframScript check for exact symbolic algebra when the claim is
algebraic.
\item A fresh adversarial verifier subagent has accepted the proof-tree node.
\end{enumerate}
\end{definition}

\begin{warning}[Theorem targets are not theorems yet]
Statements labelled ``target'' or ``conjecture'' are not claimed as proved.
They are included so that the document contains the full mathematical surface
area of the project.
\end{warning}

\section{Local Source Provenance}

\begin{longtable}{p{0.23\linewidth}p{0.67\linewidth}}
\textbf{Source identifier} & \textbf{Local source and use} \\
\hline
\texttt{SRC-FIB-TREBST-LOCAL} &
\texttt{references/FibonacciAnyonModels.pdf}; extracted text in
\texttt{references/text/FibonacciAnyonModels.txt}. Used for Fibonacci fusion
rules, quantum dimension, the standard Fibonacci \(F\)-matrix, and finite
Fibonacci \(R\)- and braid-matrix data. \\
\texttt{SRC-FIB-TREBST-ARXIV} &
\texttt{references/TrebstShortIntroductionFibonacciAnyons.pdf}; extracted
text in \texttt{references/text/TrebstShortIntroductionFibonacciAnyons.txt}.
Used as a second local source for the golden-ratio quantum dimension. \\
\texttt{SRC-GOLDEN-CHAIN} &
\texttt{references/GoldenChainFeiguinEtAl.pdf}; extracted text in
\texttt{references/text/GoldenChainFeiguinEtAl.txt}. Used for the same
Fibonacci \(F\)-matrix. \\
\texttt{SRC-KZ-FIB} &
\texttt{references/IsingLikeFibonacciAnyonsKZ.pdf}; extracted text in
\texttt{references/text/IsingLikeFibonacciAnyonsKZ.txt}. Used for the
conformal weight \(h_\tau=2/5\). The extracted text loses the fraction slash
in the table row, so the row was checked visually against the local PDF page:
the spin-1 field has \(h=2/5\). \\
\texttt{SRC-TL-JONES} &
\texttt{references/TemperleyLiebRootsJonesQuotient.pdf}; extracted text in
\texttt{references/text/TemperleyLiebRootsJonesQuotient.txt}. Used for the
Temperley--Lieb and Jones quotient background. \\
\texttt{SRC-ISING-CHAIN} &
\texttt{references/PoilblancOneDimensionalItinerantInteractingAnyons.pdf};
extracted text in
\texttt{references/text/PoilblancOneDimensionalItinerantInteractingAnyons.txt}.
Used for Ising labels, fusion rules, and the nonchiral dimensions
\(\Delta_\sigma=1/8\), \(\Delta_\psi=1\). \\
\texttt{SRC-STRING-NET} &
\texttt{references/StringNetCondMat0510613.pdf}; extracted text in
\texttt{references/text/StringNetCondMat0510613.txt}. Used for string-net
fixed-point data: \(F\)-symbols and quantum dimensions. \\
\texttt{SRC-PENNEYS-UFC} &
\texttt{references/text/PenneysUnitaryFusionCategories.md}. Used for unitary
fusion category background, direct sums, orthogonal direct sums, and polar
decomposition language. \\
\texttt{SRC-DEFECT-FOCK} &
\texttt{references/GaugingDefectsQuantumSpinSystems.pdf}; extracted text in
\texttt{references/text/GaugingDefectsQuantumSpinSystems.txt}. Used for
fusion category definitions, Fock-space direct sums, and fusion-tree Hilbert
spaces. \\
\texttt{SRC-OAR-FERMIONS} and \texttt{SRC-OAR-WAVELETS} &
Local extracted texts in \texttt{references/text/}. Used for
operator-algebraic renormalization and isometric wavelet refinement maps. \\
\end{longtable}

\section{Categorical Background}

\begin{definition}[Fusion category, local working form]
Over the complex numbers, a fusion category is treated as a rigid semisimple
tensor category with finitely many isomorphism classes of simple objects and a
simple tensor unit. The local sources are
\texttt{references/text/GaugingDefectsQuantumSpinSystems.txt:400} and
\texttt{:404}.
\end{definition}

\begin{definition}[Unitary fusion category, local working form]
A unitary fusion category is a fusion category with a dagger operation on
morphisms and a compatible unitary monoidal structure. Orthogonal direct sums
have isometric inclusions whose adjoints are projections. The local sources
are \texttt{references/text/PenneysUnitaryFusionCategories.md:219},
\texttt{:319}, \texttt{:456}, and \texttt{:458}.
\end{definition}

\begin{definition}[Skeletal finite model]
For Lean formalisation, the initial model consists of a finite type of simple
labels, a distinguished unit label, and natural-number fusion multiplicities
\[
N_{ab}^{c}\in \mathbb{N}.
\]
The local skeletal-data source states that one chooses a finite set of simple
objects with a distinguished unit and finite-dimensional fusion spaces. The
Lean file \texttt{Formalisation/Foundations/SkeletalFusion.lean} formalises
only the finite label/multiplicity skeleton with left and right unit laws:
\[
N_{\mathbf 1 a}^{c}=\delta_{a,c},\qquad
N_{a\mathbf 1}^{c}=\delta_{a,c}.
\]
It also proves that the total unit fusion multiplicity is \(1\) on either
side. AF node \texttt{1.2.7} has been accepted by verifier
\texttt{verifier-finite-skeletal-fusion-data-001}. This does not construct a
fusion category, associators, \(F\)-symbols, pentagon equations, duals,
Hilbert morphism spaces, or categorical reconstruction.
\end{definition}

\begin{definition}[Local occupation object]
Let \(\mathcal C\) be the finite skeletal model. Define
\[
Q:=\bigoplus_{a\in\Irr(\mathcal C)} a.
\]
The vacuum summand represents an empty site. For \(N\) ordered sites, define
\[
Q^{\otimes N}:=\underbrace{Q\otimes \cdots \otimes Q}_{N\text{ factors}}.
\]
Status: the categorical direct-sum object is still a target. The accepted
finite coordinate skeleton is a finite label type with a distinguished vacuum
label; the summand labels of \(Q\) are represented by that label type, and the
summand labels of \(Q^{\otimes N}\) are ordered configurations
\(\mathrm{Fin}(N)\to A\). This is formalised as
\texttt{LocalOccupationModel} and \texttt{LocalOccupationTensorSummand} in
\texttt{Formalisation/Foundations/ProjectDefinitions.lean}.
\end{definition}

\begin{definition}[Indefinite-particle Hilbert space]
For a total charge \(W\in\Irr(\mathcal C)\), define
\[
\mathcal H_N^W:=\Hom_{\mathcal C}(W,Q^{\otimes N}).
\]
If \(\mathcal C\) is only a fusion category, this is a finite-dimensional
complex vector space. If \(\mathcal C\) is unitary, it is treated as a Hilbert
space after choosing the standard inner products on fusion-tree bases.
Status: the categorical Hom-space definition remains a target. The accepted
finite coordinate skeleton assumes a finite coordinate basis for every
configuration-sector \(\Hom(W,a_1\otimes\cdots\otimes a_N)\). The flattened
coordinates of the resulting fixed sector are formalised as
\texttt{IndefiniteParticleSectorCoordinates}, with component coordinates and
pairing preservation inherited from the configuration-space coordinate
expansion.
\end{definition}

\begin{definition}[Particle number]
For a configuration \(\vec a=(a_1,\ldots,a_N)\), define its particle number by
\[
\#\{i:a_i\ne \mathbf 1\}.
\]
Thus \(\mathcal H_N^W\) contains sectors with particle number from \(0\) to
\(N\). In the finite Lean model, an ordered configuration of length \(N\) is a
function \(\mathrm{Fin}(N)\to A\), where \(A\) is the finite label type.
\end{definition}

\section{General Theorem Targets}

\begin{theorem}[Finite ordered configuration combinatorics]
Let \(A\) be a finite label set and let \(N\) be a natural number. The set of
ordered configurations \(\mathrm{Fin}(N)\to A\) has cardinality
\[
|A|^N.
\]
If a vacuum label is chosen, the particle number of any configuration is at
most \(N\). The constant vacuum configuration has particle number \(0\), and a
constant non-vacuum configuration has particle number \(N\). For Fibonacci
labels \(\{\mathbf 1,\tau\}\), the number of configurations is \(2^N\).
These finite combinatorial statements are proved in
\texttt{Formalisation/Foundations/Configurations.lean}; Julia and
WolframScript check Fibonacci configuration counts and particle sectors through
\(N=8\). AF node \texttt{1.2.1} has been accepted by verifier
\texttt{verifier-config-001}. This theorem does not prove the categorical
Hom-space expansion below.
\end{theorem}

\begin{theorem}[Finite direct-sum coordinate model]
Let \(I\) be a finite set and let \(K_i\) be a finite coordinate basis for
the \(i\)-th component. The space of flattened complex coordinates
\[
\{x:\coprod_{i\in I}K_i\to \mathbb C\}
\]
is complex-linearly equivalent to the dependent family of component coordinate
spaces
\[
\prod_{i\in I}\{x_i:K_i\to\mathbb C\}.
\]
The equivalence sends \(x\) to the family \(x_i(k)=x(i,k)\), and its inverse
sends a family \(x_i\) to the flattened function \((i,k)\mapsto x_i(k)\).
It also satisfies
\[
\left|\coprod_{i\in I}K_i\right|=\sum_{i\in I}|K_i|
\]
and preserves the standard finite coordinate pairing
\[
\sum_{i\in I}\sum_{k\in K_i}\overline{x_i(k)}\,y_i(k)
=
\sum_{(i,k)\in\coprod_i K_i}\overline{x(i,k)}\,y(i,k).
\]
This is proved in
\texttt{Formalisation/Foundations/DirectSumCoordinates.lean}; Julia and
WolframScript check exact rational complex examples. AF node \texttt{1.2.2}
has been accepted by verifier \texttt{verifier-dirsum-coord-001}. This result
is only the finite coordinate skeleton obtained after bases have been chosen
in the summands; it does not construct categorical biproducts, \(Q\), tensor
powers, Hom functors, dagger categories, or fusion-category basis provenance.
\end{theorem}

\begin{theorem}[Finite direct-sum projection equations]
In the same finite coordinate model, write
\[
\iota_i:\{x_i:K_i\to\mathbb C\}\to
\{x:\coprod_{j\in I}K_j\to\mathbb C\}
\]
for the inclusion that places component coordinates into the \(i\)-th block,
and write
\[
\pi_i:\{x:\coprod_{j\in I}K_j\to\mathbb C\}\to
\{x_i:K_i\to\mathbb C\}
\]
for the projection onto the \(i\)-th block. Then
\[
\pi_i\iota_i=\operatorname{id},
\qquad
\pi_i\iota_j=0\quad(i\ne j),
\qquad
\sum_{i\in I}\iota_i\pi_i=\operatorname{id}.
\]
This is the finite coordinate analogue of Penneys'
\(\pi_i\circ\iota_j=\delta_{i=j}\operatorname{id}\) and
\(\sum_i\iota_i\circ\pi_i=\operatorname{id}\) equations. It is proved in
\texttt{Formalisation/Foundations/DirectSumCoordinates.lean}; Julia and
WolframScript check exact finite matrix examples. AF node \texttt{1.2.3} has
been accepted by verifier \texttt{verifier-dirsum-proj-001}. This result does
not construct categorical biproducts or the direct-sum universal property.
\end{theorem}

\begin{theorem}[Finite orthogonal direct-sum coordinates]
In the finite coordinate model, equip every component and the flattened direct
sum with the standard complex coordinate pairing. Then each component
inclusion \(\iota_i\) is isometric:
\[
\langle \iota_i x,\iota_i y\rangle
=
\langle x,y\rangle .
\]
The coordinate range map \(P_i=\iota_i\pi_i\) satisfies
\[
P_i^2=P_i,
\qquad
P_iP_j=0\quad(i\ne j).
\]
Moreover, component projection is the coordinate adjoint of component
inclusion:
\[
\langle \iota_i x,y\rangle
=
\langle x,\pi_i y\rangle,
\qquad
\langle y,\iota_i x\rangle
=
\langle \pi_i y,x\rangle,
\]
and therefore the coordinate range projection is self-adjoint for the
pairing:
\[
\langle P_i x,y\rangle=\langle x,P_i y\rangle .
\]
This is the finite coordinate analogue of Penneys' statement that orthogonal
direct-sum inclusions are isometries and the maps
\(\iota_i\iota_i^\dagger\) are mutually orthogonal projections. It is proved
in \texttt{Formalisation/Foundations/DirectSumCoordinates.lean}; Julia and
WolframScript check exact finite matrix and pairing examples. AF nodes
\texttt{1.2.4} and \texttt{1.2.8} have been accepted by verifiers
\texttt{verifier-dirsum-orth-001} and
\texttt{verifier-dirsum-adjoint-coord-001}. This result does not construct
categorical adjoints or orthogonal direct sums.
\end{theorem}

\begin{theorem}[Finite configuration-space coordinate expansion]
Let \(A\) be a finite label set and let
\[
\operatorname{Config}(A,N)=\{\,c:\operatorname{Fin}(N)\to A\,\}
\]
be the finite set of ordered configurations. Suppose that every
configuration \(c\) has been assigned a finite coordinate basis \(B_c\). The
space of flattened complex coordinates
\[
\{x:\coprod_{c\in\operatorname{Config}(A,N)}B_c\to\mathbb C\}
\]
is complex-linearly equivalent to the dependent family of component
coordinate spaces
\[
\prod_{c\in\operatorname{Config}(A,N)}\{x_c:B_c\to\mathbb C\}.
\]
The equivalence sends \(x\) to the family \(x_c(b)=x(c,b)\). Its coordinate
cardinality satisfies
\[
\left|\coprod_{c\in\operatorname{Config}(A,N)}B_c\right|
=\sum_{c\in\operatorname{Config}(A,N)}|B_c|.
\]
If all sectors have the same finite basis \(B\), then
\[
\left|\coprod_{c\in\operatorname{Config}(A,N)}B\right|
=|A|^N\,|B|.
\]
The equivalence preserves the standard finite coordinate pairing:
\[
\sum_c\sum_{b\in B_c}\overline{x_c(b)}\,y_c(b)
=
\sum_{(c,b)}\overline{x(c,b)}\,y(c,b).
\]
For Fibonacci labels \(A=\{\mathbf 1,\tau\}\), the constant-sector
cardinality is \(2^N|B|\). Lean proves these statements in
\texttt{Formalisation/Foundations/ConfigurationSpace.lean} using the accepted
configuration and direct-sum coordinate results. Julia and WolframScript
check exact Fibonacci \(N=3\) examples with dependent sector sizes and
rational complex pairings. AF node \texttt{1.2.5} has been accepted by
verifier \texttt{verifier-config-space-coord-001}. This theorem is only the
finite, basis-dependent coordinate skeleton after sector bases have been
supplied; it does not construct categorical Hom spaces, \(Q^{\otimes N}\),
tensor products, canonical Hom isomorphisms, categorical direct sums, or
unitarity in a fusion category.
\end{theorem}

\begin{theorem}[Finite coordinate skeleton of \(Q\) and \(\mathcal H_N^W\)]
Let \(A\) be a finite label type with a distinguished vacuum label. In the
finite coordinate model, the summands of \(Q\) are indexed by \(A\), and the
summands of \(Q^{\otimes N}\) are indexed by functions
\[
\mathrm{Fin}(N)\to A.
\]
Hence there are \(|A|^N\) label summands in \(Q^{\otimes N}\), and the constant
vacuum configuration has particle number \(0\). After choosing a finite basis
for each sector \(\Hom(W,a_1\otimes\cdots\otimes a_N)\), the fixed sector
\(\mathcal H_N^W\) has flattened coordinates over all configuration-sector
basis vectors. These coordinates expand to component coordinates over
configurations, and the standard finite coordinate pairing is preserved.

Lean proves these statements in
\texttt{Formalisation/Foundations/ProjectDefinitions.lean}. Julia and
WolframScript check the Fibonacci two-label instance, the \(Q^{\otimes 4}\)
label count, vacuum and all-\(\tau\) particle numbers, flattened coordinate
dimension, unflattening, and pairing preservation. AF node \texttt{1.2.6} has
been accepted by verifier \texttt{verifier-project-q-hilb-coord-001}. This is
only a finite, basis-dependent coordinate skeleton; it does not construct
categorical direct sums, tensor powers, Hom functors, canonical categorical
isomorphisms, or fusion-category Hilbert-space semantics.
\end{theorem}

\begin{theorem}[Finite truncated distinguishable-Fock coordinates]
Let \(B\) be a finite local coordinate basis and let \(M\) be a nonnegative
integer. The finite truncation of distinguishable Fock coordinates through
particle number \(M\) has basis
\[
\coprod_{n=0}^{M}\{\,w:\operatorname{Fin}(n)\to B\,\}.
\]
Equivalently, a basis vector is a particle number \(n\le M\) together with an
ordered word of length \(n\) in the local basis. Flattened coordinates on this
coproduct are complex-linearly equivalent to the family of fixed-particle
number coordinate vectors. The basis cardinality is
\[
\left|\coprod_{n=0}^{M}\{w:\operatorname{Fin}(n)\to B\}\right|
=\sum_{n=0}^{M}|B|^n.
\]
The zero-particle sector has one basis vector, and the one-particle sector has
\(|B|\) basis vectors. The flattening by particle number preserves the
standard finite coordinate pairing. For the Fibonacci two-label local basis,
the cardinality is \(\sum_{n=0}^{M}2^n\).

This is the finite coordinate truncation of the locally sourced
distinguishable-Fock construction as a direct sum over particle number, with a
one-dimensional zero-particle vacuum sector. Lean proves the finite statements
in \texttt{Formalisation/Foundations/FockSpace.lean}. Julia and WolframScript
check exact finite counts and rational complex pairings. AF node
\texttt{1.2.9} has been accepted by verifier
\texttt{verifier-truncated-fock-coord-001}. This result does not construct an
infinite direct sum, a Hilbert-space completion, tensor-product Hilbert
spaces, exchange-symmetry quotients, categorical Hom spaces, or physical
Fock-space semantics.
\end{theorem}

\begin{theorem}[Configuration-space expansion target]
There should be a canonical isomorphism
\[
\mathcal H_N^W\cong
\bigoplus_{(a_1,\ldots,a_N)}
\Hom_{\mathcal C}(W,a_1\otimes\cdots\otimes a_N).
\]
If the category is unitary and the direct-sum inclusions are isometries, this
isomorphism should be unitary. Status: pending. The finite configuration-space
coordinate theorem above proves only the basis-dependent vector-space
skeleton of this target after sector bases have been supplied.
\end{theorem}

\begin{theorem}[Finite Fibonacci braid matrix relation]
Let
\[
B_{12}=\begin{pmatrix}q^2&0\\0&-q\end{pmatrix},
\qquad
B_{23}=c\begin{pmatrix}-1&q s\\q s&q^2\end{pmatrix}.
\]
If
\[
c(q+1)=q,\qquad qs^2=q^2+q+1,
\]
then
\[
B_{12}B_{23}B_{12}=B_{23}B_{12}B_{23}.
\]
This is the algebraic core of the sourced Fibonacci relation with
\[
q=\exp(-2\pi i/5),\qquad
\phi=(1+\sqrt5)/2,\qquad
s=\sqrt{\phi},\qquad
c=q/(q+1),
\]
where the source writes \(B_{23}=q(q+1)^{-1}
\begin{pmatrix}-1&q\sqrt\phi\\q\sqrt\phi&q^2\end{pmatrix}\) and
\(B_{12}B_{23}B_{12}-B_{23}B_{12}B_{23}=0\).
Lean proves the parameterised algebra in
\texttt{Formalisation/Fibonacci/BraidMatrices.lean}; Julia checks the
specialisation numerically at high precision; WolframScript checks it exactly.
AF node \texttt{1.3.10} has been accepted by verifier
\texttt{verifier-fib-braid-001}. This result is not a categorical braid group
representation theorem.
\end{theorem}

\begin{theorem}[Finite braid-matrix unitarity by conjugation]
Let \(I\) be a finite index set and let \(F,R\in M_I(\mathbb C)\). Assume
\[
F^\dagger F=I,\qquad FF^\dagger=I,\qquad
R^\dagger R=I,\qquad RR^\dagger=I.
\]
Define
\[
B=F R F^\dagger .
\]
Then
\[
B^\dagger B=I,\qquad BB^\dagger=I.
\]
This is the finite matrix algebra behind the locally sourced statements that
fusion-tree changes are represented by unitary \(F\)-matrices, braidings give
unitary \(R\)-matrices that are diagonal in a suitable basis, and a braid
matrix is
\[
B=F R F^{-1}.
\]
For the sourced five-dimensional Fibonacci adjacent-label basis, the matrix
\(F\) has three \(1\)'s and the nontrivial Fibonacci block
\[
\begin{pmatrix}\phi^{-1}&\phi^{-1/2}\\
\phi^{-1/2}&-\phi^{-1}\end{pmatrix},
\]
and
\[
R=\operatorname{diag}\bigl(e^{4\pi i/5},e^{-3\pi i/5},e^{-3\pi i/5},
e^{4\pi i/5},e^{-3\pi i/5}\bigr).
\]
Lean proves the generic finite-matrix theorem in
\texttt{Formalisation/LinearAlgebra/Isometry.lean}. Julia checks this
five-dimensional sourced specialisation numerically at high precision, and
WolframScript checks it exactly. AF node \texttt{1.3.11} has been accepted by
verifier \texttt{verifier-fib-braid-unitary-001}. The extracted text around
one displayed formula is OCR-ambiguous, so the local PDF page and the later
extracted line giving \(B=FRF^{-1}\) were both used. This theorem does not
construct categorical braidings, conformal blocks, mapping-class-group
representations, or analytic monodromy.
\end{theorem}

\begin{theorem}[Braid group action target]
If \(\mathcal C\) is braided, the neighbouring exchange maps induced by
\[
c_{Q,Q}:Q\otimes Q\to Q\otimes Q
\]
should satisfy
\[
B_iB_{i+1}B_i=B_{i+1}B_iB_{i+1},\qquad
B_iB_j=B_jB_i\quad |i-j|\geq 2.
\]
Status: pending.
\end{theorem}

\begin{definition}[Fine-graining isometry]
For \(M>N\), a fine-graining map is a morphism or finite matrix
\[
E_{N\to M}:Q^{\otimes N}\to Q^{\otimes M}
\]
such that
\[
E_{N\to M}^{\dagger}E_{N\to M}=\operatorname{id}_{Q^{\otimes N}}.
\]
For a charge sector \(W\), the induced map is
\[
V_{N\to M}^W(f)=E_{N\to M}\circ f.
\]
In finite coordinate form, after bases have been chosen, this definition is
formalised as a structure consisting of a matrix \(E\) and the equation
\[
E^\dagger E=I.
\]
The Lean structure is \texttt{FiniteFineGraining} in
\texttt{Formalisation/LinearAlgebra/FineGraining.lean}. It packages the
validated finite consequences: ascending unitality, postcomposition Gram
preservation, state-embedding trace preservation, expectation preservation,
ascending adjoint preservation
\[
E^\dagger O^\dagger E=(E^\dagger O E)^\dagger,
\]
logical-lift retraction, logical lift of the identity
\[
(E I)E^\dagger=EE^\dagger,
\]
and the code projection identities for \(EE^\dagger\). AF node \texttt{1.4.12}
has been accepted by verifier
\texttt{verifier-finite-fine-graining-def-001}; AF node \texttt{1.4.14} has
been accepted by verifier \texttt{verifier-finite-fg-star-lift-001} for the
packaged adjoint-preservation and logical-identity consequences. This is only
the finite coordinate definition; it does not construct \(Q^{\otimes N}\),
categorical morphisms, Hom-sector semantics, categorical channel semantics, or
physical fine-graining maps in a fusion category.
\end{definition}

\begin{theorem}[Postcomposition isometry target]
If \(E^\dagger E=\operatorname{id}\), then
\(V^W(f)=E\circ f\) is an isometry for every \(W\). Conversely, if all
simple-sector maps \(V^W\) are isometries and the Hom functor separates
morphisms, then \(E^\dagger E=\operatorname{id}\). Status: the finite-matrix
forward direction is proved in
\texttt{Formalisation/LinearAlgebra/Postcompose.lean}: if \(E^\dagger E=I\),
then
\[
(EF)^\dagger(EG)=F^\dagger G.
\]
AF node \texttt{1.4.4} is accepted for this coordinate matrix theorem.
The finite square-test converse is also proved: if
\[
(EF)^\dagger(EG)=F^\dagger G
\]
for every pair of square coarse-coordinate matrices \(F,G\), then
\[
E^\dagger E=I.
\]
Indeed, take \(F=I\) and \(G=I\). AF node \texttt{1.4.8} is accepted for
this finite matrix converse. The full categorical converse using
semisimplicity and a separating Hom functor, and the categorical Hom-sector
semantics, remain pending.
\end{theorem}

\begin{theorem}[Component orthogonality target]
For decompositions
\[
Q^{\otimes N}\cong \bigoplus_{\vec a}A_{\vec a},
\qquad
Q^{\otimes M}\cong \bigoplus_{\vec b}B_{\vec b},
\]
define
\[
E_{\vec b,\vec a}:=\pi_{\vec b}E\iota_{\vec a}.
\]
Then \(E\) is an isometry exactly when
\[
\sum_{\vec b}E_{\vec b,\vec a}^{\dagger}E_{\vec b,\vec a'}
=\delta_{\vec a,\vec a'}\operatorname{id}_{A_{\vec a}}.
\]
Status: the scalar finite-matrix version is proved in
\texttt{Formalisation/LinearAlgebra/Component.lean}: for a matrix \(E\),
\[
E^\dagger E=I
\quad\Longleftrightarrow\quad
\sum_b \overline{E_{b,a}}E_{b,a'}=\delta_{a,a'}.
\]
AF node \texttt{1.4.2} is accepted only for this coordinate matrix theorem.
Block Hom spaces and categorical direct-sum semantics remain pending.
\end{theorem}

\begin{theorem}[Block tensor isometry target]
Let \(M>N\), let the fine sites be partitioned into ordered blocks of sizes
\(r_i\geq 1\), and let
\[
\eta_i:Q\to Q^{\otimes r_i}
\]
be isometries. Then
\[
E_{N\to M}^{\mathrm{block}}=\eta_1\otimes\cdots\otimes\eta_N
\]
is an isometry. Status: the two-factor finite-matrix version is proved in
\texttt{Formalisation/LinearAlgebra/Tensor.lean}: if \(E_1^\dagger E_1=I\)
and \(E_2^\dagger E_2=I\), then
\[
(E_1\otimes E_2)^\dagger(E_1\otimes E_2)=I.
\]
AF node \texttt{1.4.3} is accepted only for this Kronecker product theorem.
The same Lean file also proves the heterogeneous three-factor finite matrix
case: if \(E_1^\dagger E_1=I\), \(E_2^\dagger E_2=I\), and
\(E_3^\dagger E_3=I\), then
\[
((E_1\otimes E_2)\otimes E_3)^\dagger
((E_1\otimes E_2)\otimes E_3)=I.
\]
AF node \texttt{1.4.10} has been accepted by verifier
\texttt{verifier-tensor-three-het-001}.
The fixed-local-map induction is proved in
\texttt{Formalisation/LinearAlgebra/TensorPower.lean}: if \(E^\dagger E=I\),
then every finite Kronecker tensor power \(E^{\otimes N}\) is an isometry.
Julia and WolframScript check exact rational examples through \(N=6\).
AF node \texttt{1.4.7} has been accepted by verifier
\texttt{verifier-tensor-power-001}.
The arbitrary heterogeneous finite-coordinate induction is proved in
\texttt{Formalisation/LinearAlgebra/HeterogeneousTensor.lean}. For any
natural number \(N\), finite dependent families of fine and coarse coordinate
index types, and local matrices \(E_i\) with \(E_i^\dagger E_i=I\), the
recursively ordered heterogeneous Kronecker product satisfies
\[
\bigl(E_0\otimes E_1\otimes\cdots\otimes E_{N-1}\bigr)^\dagger
\bigl(E_0\otimes E_1\otimes\cdots\otimes E_{N-1}\bigr)=I.
\]
AF node \texttt{1.4.16} has been accepted by verifier
\texttt{verifier-het-tensor-n-001}. Categorical tensor products, canonical
associators, ordered partition semantics, and construction of local maps from
fusion-category data remain pending.
\end{theorem}

\begin{theorem}[Unitary dressing target]
If \(U_M:Q^{\otimes M}\to Q^{\otimes M}\) is unitary and each \(\eta_i\) is
isometric, then
\[
E_{N\to M}=U_M(\eta_1\otimes\cdots\otimes\eta_N)
\]
is isometric. Status: the finite-matrix calculation
\((UE)^\dagger(UE)=I\) from \(U^\dagger U=I\) and \(E^\dagger E=I\) is proved
in \texttt{Formalisation/LinearAlgebra/Isometry.lean}. The node has been
accepted by verifier \texttt{verifier-mat-dress-001} only for this
finite-matrix calculation. For three heterogeneous factors, the Lean file
\texttt{Formalisation/LinearAlgebra/Tensor.lean} also proves that if
\(U\) is unitary on the three-factor fine product, then
\[
\bigl(U((E_1\otimes E_2)\otimes E_3)\bigr)^\dagger
U((E_1\otimes E_2)\otimes E_3)=I.
\]
AF node \texttt{1.4.11} has been accepted by verifier
\texttt{verifier-dressed-three-tensor-001}.
For fixed-local-map tensor powers, the Lean file
\texttt{Formalisation/LinearAlgebra/TensorPower.lean} proves that if
\(E^\dagger E=I\), \(E^{\otimes N}\) is the recursively defined Kronecker
tensor power, and \(U\) is unitary on the fine tensor-power coordinate space,
then
\[
\bigl(U E^{\otimes N}\bigr)^\dagger\bigl(U E^{\otimes N}\bigr)=I.
\]
Julia and WolframScript check exact rational rectangular examples with a
nontrivial reversal permutation unitary through \(N=4\). AF node
\texttt{1.4.15} has been accepted by verifier
\texttt{verifier-dressed-tensor-power-001}.
The heterogeneous finite-coordinate theorem in
\texttt{Formalisation/LinearAlgebra/HeterogeneousTensor.lean} also proves the
arbitrary \(N\)-factor dressed case over the ordered recursive tensor index:
if \(U\) is unitary on the fine tensor-product coordinate space, then
\[
\bigl(U(E_0\otimes\cdots\otimes E_{N-1})\bigr)^\dagger
U(E_0\otimes\cdots\otimes E_{N-1})=I.
\]
This is part of AF node \texttt{1.4.16}. Full categorical unitary dressing,
ordered partition semantics, and physical construction of \(U\) remain
pending.
\end{theorem}

\section{Finite-Dimensional Linear Algebra}

\begin{theorem}[Polar section target]
Let \(B:H_{\mathrm{fine}}\to H_{\mathrm{coarse}}\) be a full-rank linear map
between finite-dimensional Hilbert spaces. Define
\[
E:=B^\dagger(BB^\dagger)^{-1/2}.
\]
Then \(E^\dagger E=\operatorname{id}_{H_{\mathrm{coarse}}}\). Status: the
conditional finite-matrix algebra is proved in
\texttt{Formalisation/LinearAlgebra/Polar.lean}. If a matrix \(R\) satisfies
\[
R^\dagger(BB^\dagger)R=I,
\]
then \(E=B^\dagger R\) satisfies \(E^\dagger E=I\). AF node \texttt{1.4.5}
is accepted only for this conditional statement; existence and positivity of
\((BB^\dagger)^{-1/2}\) remain pending.

There is also an accepted diagonal special case. Suppose \(BB^\dagger\) is
diagonal in the chosen coarse basis, with diagonal entries \(g_\mu>0\). Define
\[
R_{\mu\mu'}=
\begin{cases}
(\sqrt{g_\mu})^{-1}, & \mu=\mu',\\
0, & \mu\ne\mu'.
\end{cases}
\]
Then \(R^\dagger=R\),
\[
R^\dagger(BB^\dagger)R=I,
\]
and \(B^\dagger R\) is an isometry. Lean defines this matrix as
\texttt{diagonalInvSqrt} and proves the inverse-square-root condition and
isometry in \texttt{Formalisation/LinearAlgebra/Polar.lean}. Julia checks the
diagonal case \(g=(4,9)\), and WolframScript checks the same case exactly. AF
node \texttt{1.4.13} has been accepted by verifier
\texttt{verifier-polar-diag-sqrt-001}. This does not prove the general
positive matrix square-root theorem, full polar decomposition existence, the
full-rank criterion, or categorical Hilbert-space semantics.
\end{theorem}

\begin{theorem}[Amplitude formula target]
For orthonormal bases, write
\[
B|\nu\rangle_{\mathrm{fine}}
=\sum_\mu \beta_{\mu\nu}|\mu\rangle_{\mathrm{coarse}}
\]
and
\[
G_{\mu\mu'}=\sum_\nu
\beta_{\mu\nu}\overline{\beta_{\mu'\nu}}.
\]
Then \(E|\mu\rangle=\sum_\nu A_{\nu\mu}|\nu\rangle\) with
\[
A_{\nu\mu}=\sum_{\mu'}
\overline{\beta_{\mu'\nu}}\,(G^{-1/2})_{\mu'\mu}.
\]
If \(\mu\) does not mix with other coarse labels, then
\[
A_\nu^\mu=
\frac{\overline{\beta_{\mu\nu}}}
{\sqrt{\sum_{\nu'}|\beta_{\mu\nu'}|^2}}.
\]
Status: \texttt{Formalisation/LinearAlgebra/Polar.lean} proves the coordinate
entry formula for \(E=B^\dagger R\), with \(R\) standing for the
inverse-square-root factor. The finite one-row no-mixing corollary is proved
in \texttt{Formalisation/LinearAlgebra/NoMixing.lean}: for a blocking row
\(\beta_\nu\), the matrix \(B^\dagger R\) has entries
\[
\frac{\overline{\beta_\nu}}
{\sqrt{\sum_{\nu'}|\beta_{\nu'}|^2}}
\]
and is an isometry when the denominator is nonzero. AF node \texttt{1.4.9} is
accepted for exactly this finite matrix corollary. Full square-root and polar
decomposition semantics remain pending.
\end{theorem}

\begin{theorem}[Ascending channel target]
For an isometry \(V\), the observable map
\[
\mathsf A(O)=V^\dagger O V
\]
is unital, completely positive, star-preserving, and expectation-value
preserving against embedded states. The finite matrix unital, star-preserving,
and retraction identities are formalised in
\texttt{Formalisation/LinearAlgebra/Isometry.lean}. The finite matrix trace
identities
\[
\operatorname{tr}(E\rho E^\dagger)=\operatorname{tr}(\rho)
\quad\text{when }E^\dagger E=\mathbf 1
\]
and
\[
\operatorname{tr}\bigl((E\rho E^\dagger)O\bigr)
=\operatorname{tr}\bigl(\rho(E^\dagger O E)\bigr)
\]
are formalised in \texttt{Formalisation/LinearAlgebra/Trace.lean}. The same
file proves the square-form identity
\[
K^\dagger(X^\dagger X)K=(XK)^\dagger(XK),
\]
which is a finite-dimensional positivity witness. More generally, call a
finite matrix \(O\) Gram-form positive if there is a finite witness index set
and a matrix \(X\), possibly rectangular, with
\[
O=X^\dagger X.
\]
If \(O\) is a fine observable with such a witness, then
\[
E^\dagger O E=(XE)^\dagger(XE),
\]
so the ascending congruence preserves this Gram-form witness. In Mathlib's
finite quadratic-form definition of positive semidefinite matrices, the same
file proves the order-theoretic implication
\[
O\geq 0\quad\Longrightarrow\quad E^\dagger O E\geq 0.
\]
This is the ordinary finite matrix fact that a positive semidefinite matrix
remains positive semidefinite after congruence by \(E\). For every finite
auxiliary coordinate set \(R\), the same file also proves the amplified
finite statement
\[
O\geq 0\quad\Longrightarrow\quad
(I_R\otimes E)^\dagger O(I_R\otimes E)\geq 0
\]
for matrices \(O\) on \(R\times\text{fine}\). This proves amplified
positivity for the one-Kraus congruence \(I_R\otimes E\). If
\(\{K_i\}_{i\in S}\) is any finite family of matrices with common source and
target, then the finite congruence sums also preserve positivity:
\[
O\geq 0\quad\Longrightarrow\quad
\sum_{i\in S}K_i^\dagger O K_i\geq 0
\]
and, for every finite auxiliary coordinate set \(R\),
\[
O\geq 0\quad\Longrightarrow\quad
\sum_{i\in S}(I_R\otimes K_i)^\dagger O(I_R\otimes K_i)\geq 0.
\]
For a finite family of matrices \(\{K_i\}_{i\in S}\), define
\[
\mathsf A(O)=\sum_{i\in S}K_i^\dagger O K_i,\qquad
\mathsf S(\rho)=\sum_{i\in S}K_i\rho K_i^\dagger.
\]
If
\[
\sum_{i\in S}K_i^\dagger K_i=I,
\]
then
\[
\mathsf A(I)=I,\qquad
\operatorname{tr}(\mathsf S(\rho))=\operatorname{tr}(\rho).
\]
For every finite family, the same finite matrix formalisation proves
\[
\mathsf A(O^\dagger)=\mathsf A(O)^\dagger
\]
and
\[
\operatorname{tr}(\mathsf S(\rho)O)
=\operatorname{tr}(\rho\,\mathsf A(O)).
\]
This is finite-dimensional matrix algebra. It does not assert that a
particular physical refining channel has a sourced Kraus representation. For
the logical lift
\[
\mathsf L(O)=VOV^\dagger,
\]
the finite matrix formalisation also proves
\[
\mathsf L(I)=VV^\dagger,\qquad
(VV^\dagger)^2=VV^\dagger,\qquad
(VV^\dagger)^\dagger=VV^\dagger.
\]
Thus the lifted coarse identity is the code-subspace projection in the finite
matrix sense, not the full fine identity. Status: AF nodes \texttt{1.5.1},
\texttt{1.5.2}, \texttt{1.5.8}, \texttt{1.5.9}, \texttt{1.5.11}, and
\texttt{1.5.12}, \texttt{1.5.13}, and \texttt{1.5.14} are accepted for these finite-matrix
identities only. Full categorical channel semantics, categorical subobject
semantics, full spectral projection theory, and a sourced physical
Kraus-representation theorem remain pending.
\end{theorem}

\section{Fibonacci Data}

\begin{definition}[Fibonacci labels and fusion rules]
The Fibonacci model has two simple labels, \(\mathbf 1\) and \(\tau\), with
\[
\mathbf 1\otimes\mathbf 1=\mathbf 1,\quad
\mathbf 1\otimes\tau=\tau,\quad
\tau\otimes\mathbf 1=\tau,\quad
\tau\otimes\tau=\mathbf 1\oplus\tau.
\]
Local source: \texttt{references/text/FibonacciAnyonModels.txt:128},
\texttt{:129}, and \texttt{:133}. Status: formalised in
\texttt{Formalisation/Fibonacci/FusionRules.lean}. AF node \texttt{1.3.8}
has been accepted by verifier \texttt{verifier-fib-fusion-001} only for the
finite multiplicity table. The same file instantiates the finite skeletal
fusion-data structure from AF node \texttt{1.2.7} and proves the
multiplicity-free condition and unit total multiplicities for this Fibonacci
table.
\end{definition}

\begin{definition}[Golden ratio and total quantum dimension]
Define
\[
\varphi:=\frac{1+\sqrt 5}{2},\qquad d_{\mathbf 1}=1,\qquad d_\tau=\varphi.
\]
Define
\[
D^2:=1+\varphi^2.
\]
Local sources: \texttt{references/text/FibonacciAnyonModels.txt:216},
\texttt{:218}, \texttt{:304}, and
\texttt{references/text/TrebstShortIntroductionFibonacciAnyons.txt:217},
\texttt{:325}.
\end{definition}

\begin{theorem}[Fibonacci constant identities]
The following identities hold:
\[
\varphi^2=\varphi+1,\qquad
\varphi^{-1}=\varphi-1,\qquad
\varphi^{-2}+\varphi^{-1}=1,\qquad
D^2=2+\varphi.
\]
Status: Adversarial Proof Framework node \texttt{1.3.1} submitted for
verification and accepted by verifier \texttt{verifier-fib-alg-001}. Lean,
Julia, and WolframScript checks pass. The accepted scope is the real-number
algebra claim about \(\varphi\) and \(D^2\), not a broader category-theoretic
Fibonacci theorem.
\end{theorem}

\begin{proof}
The Lean proof is in \texttt{Formalisation/Fibonacci/Basic.lean}. It unfolds
\(\varphi\), uses \((\sqrt 5)^2=5\), and solves the real algebra. Julia checks
the equations numerically at 256-bit precision. WolframScript checks the same
equations by exact symbolic simplification.
\end{proof}

\begin{definition}[Fibonacci \(F\)-matrix]
The nontrivial \(F\)-matrix used initially is
\[
F^{\tau\tau\tau}_\tau=
\begin{pmatrix}
\varphi^{-1} & \varphi^{-1/2}\\
\varphi^{-1/2} & -\varphi^{-1}
\end{pmatrix}.
\]
Local sources: \texttt{references/text/FibonacciAnyonModels.txt:272},
\texttt{:297}, \texttt{:301}, \texttt{:304}, and
\texttt{references/text/GoldenChainFeiguinEtAl.txt:109}.
\end{definition}

\begin{theorem}[Fibonacci \(F\)-matrix exact identities]
The matrix above satisfies
\[
F^\dagger F=I,\qquad F^2=I.
\]
Status: ready for adversarial verification in node \texttt{1.3.2}; Lean,
Julia, and WolframScript checks pass for the parameterised theorem
\(\begin{psmallmatrix}a&b\\ b&-a\end{psmallmatrix}\) under \(a^2+b^2=1\).
The node has been accepted by verifier \texttt{verifier-fib-mat-001} as a
real matrix algebra claim. The categorical meaning of the \(F\)-move is not
included in that acceptance.
\end{theorem}

\begin{definition}[General binary Fibonacci map]
Let \(Q=\mathbf 1\oplus\tau\). A binary map
\[
\eta:Q\to Q\otimes Q
\]
has coefficients
\[
\eta_{\mathbf 1}=xv^{\mathbf 1}_{\mathbf 1\mathbf 1}
+yv^{\mathbf 1}_{\tau\tau}
\]
and
\[
\eta_\tau=pv^\tau_{\tau\mathbf 1}
+qv^\tau_{\mathbf 1\tau}
+rv^\tau_{\tau\tau}.
\]
\end{definition}

\begin{theorem}[Binary Fibonacci isometry target]
The map \(\eta\) is an isometry exactly when
\[
|x|^2+|y|^2=1,\qquad |p|^2+|q|^2+|r|^2=1.
\]
There is no cross-condition between the \(\mathbf 1\) and \(\tau\) sectors
because \(\Hom(\mathbf 1,\tau)=0\). Status: ready for adversarial
verification in node \texttt{1.3.6} for the coordinate matrix theorem after
choosing the orthogonal binary fusion-tree basis. The Lean proof is in
\texttt{Formalisation/Fibonacci/Binary.lean}. The node has been accepted by
verifier \texttt{verifier-fib-binary-001} only at this coordinate-matrix
level.
\end{theorem}

\begin{definition}[Perron--Frobenius binary amplitudes]
For a finite fusion channel \(a\to b\otimes c\), the project handoff proposes
the topological one-step coefficient
\[
A^a_{bc}=\sqrt{\frac{d_b d_c}{d_aD^2}},
\]
where \(d_a\) is the positive Frobenius--Perron dimension and
\[
D^2=\sum_x d_x^2.
\]
For Fibonacci labels \(d_{\mathbf 1}=1\), \(d_\tau=\varphi\), and
\(D^2=1+\varphi^2\). On the five allowed binary channels this gives
\[
\eta_{\mathbf 1}^{\mathrm{PF}}=
\frac1D v^{\mathbf 1}_{\mathbf 1\mathbf 1}
+\frac{\varphi}{D}v^{\mathbf 1}_{\tau\tau}
\]
and
\[
\eta_\tau^{\mathrm{PF}}=
\frac1D v^\tau_{\tau\mathbf 1}
+\frac1D v^\tau_{\mathbf 1\tau}
+\frac{\sqrt\varphi}{D}v^\tau_{\tau\tau}.
\]
\end{definition}

\begin{theorem}[Perron--Frobenius normalisation]
In the finite skeletal Fibonacci model, the dimension vector
\[
d_{\mathbf 1}=1,\qquad d_\tau=\varphi
\]
is an eigenvector of the \(\tau\)-fusion matrix with eigenvalue \(\varphi\).
The five coefficient equalities from
\(A^a_{bc}=\sqrt{d_b d_c/(d_aD^2)}\) are
\[
A^{\mathbf 1}_{\mathbf 1\mathbf 1}=\frac1D,\quad
A^{\mathbf 1}_{\tau\tau}=\frac{\varphi}{D},\quad
A^\tau_{\tau\mathbf 1}=\frac1D,\quad
A^\tau_{\mathbf 1\tau}=\frac1D,\quad
A^\tau_{\tau\tau}=\frac{\sqrt\varphi}{D}.
\]
The scalar normalisations are
\[
\frac{1}{D^2}+\frac{\varphi^2}{D^2}=1,\qquad
\frac1{D^2}+\frac1{D^2}+\frac{\varphi}{D^2}=1.
\]
Combined with the binary coordinate isometry theorem, these identities prove
that the \(5\times 2\) coordinate matrix for
\(\eta^{\mathrm{PF}}\) satisfies
\[
(\eta^{\mathrm{PF}})^\dagger\eta^{\mathrm{PF}}=I.
\]
The scalar identities are accepted in AF node \texttt{1.3.3}. The
coordinate-level matrix isometry is accepted in AF node \texttt{1.3.13} by
verifier \texttt{verifier-fib-pf-binary-001}. The finite Fibonacci
dimension-formula coefficient calculation is accepted in AF node
\texttt{1.3.15} by verifier
\texttt{verifier-fib-pf-dim-amp-001}. This does not prove the arbitrary
fusion-category amplitude theorem, categorical Perron--Frobenius naturality,
string-net fixed-point semantics, or construction of the fusion-tree basis.
\end{theorem}

\begin{definition}[Coassociative candidate]
The positive real candidate is
\[
\eta_{\mathbf 1}^{\mathrm{coassoc}}=
\varphi^{-1}v^{\mathbf 1}_{\mathbf 1\mathbf 1}
+\varphi^{-1/2}v^{\mathbf 1}_{\tau\tau}
\]
and
\[
\eta_\tau^{\mathrm{coassoc}}=
\varphi^{-1}v^\tau_{\tau\mathbf 1}
+\varphi^{-1}v^\tau_{\mathbf 1\tau}
+\varphi^{-3/2}v^\tau_{\tau\tau}.
\]
Status: scalar equations pass Julia and WolframScript; the full associator
proof is pending. The scalar equations and displayed two-component
\(F\)-comparison vector are proved in
\texttt{Formalisation/Fibonacci/Coassoc.lean}; this is ready for adversarial
verification in node \texttt{1.3.4}. The node has been accepted by verifier
\texttt{verifier-coassoc-scalar-001} only at this scalar level. It does not
prove categorical coassociativity, uniqueness, or associator semantics.
\end{definition}

\begin{theorem}[Positive scalar coassociative solution]
Let \(x,y,r\) be positive real numbers. Suppose
\[
x^2+y^2=1,\qquad 2x^2+r^2=1,
\]
and
\[
r^2=\varphi^{-3/2}xy.
\]
Then
\[
x=\varphi^{-1},\qquad y=\varphi^{-1/2},\qquad r=\varphi^{-3/2}.
\]
In Lean the coefficient \(\varphi^{-3/2}\) is represented as
\(\varphi^{-1}\varphi^{-1/2}\), and \(r=\varphi^{-3/2}\) is represented by
\(r=\varphi^{-1}\varphi^{-1/2}\). Lean now proves the scalar identities
\[
\varphi^{-1/2}=(\sqrt\varphi)^{-1},\qquad
\varphi^{-3/2}=\varphi^{-1}(\sqrt\varphi)^{-1},
\]
and the squared positive-root form
\[
r^2=(\varphi^{-1})^3,\qquad r=\sqrt{(\varphi^{-1})^3}
\]
for the positive candidate \(r\). The proof is in
\texttt{Formalisation/Fibonacci/Coassoc.lean}; Julia checks the eliminated
polynomial roots and fractional-power identities numerically at high
precision; WolframScript checks the exact quantified uniqueness implication
and fractional-power identities. AF node \texttt{1.3.12} has been accepted by
verifier \texttt{verifier-coassoc-unique-001}; AF node \texttt{1.3.16} has
been accepted by verifier \texttt{verifier-coassoc-frac-pow-001}. This
remains only scalar algebra and does not prove categorical coassociativity,
associator semantics, fusion-tree basis comparison, or existence of a
categorical map.
\end{theorem}

\begin{theorem}[Coordinate isometry of the strict binary candidate]
Use the fixed binary coordinate basis whose rows are the five allowed
channels
\[
\mathbf 1\to\mathbf 1\mathbf 1,\quad
\mathbf 1\to\tau\tau,\quad
\tau\to\tau\mathbf 1,\quad
\tau\to\mathbf 1\tau,\quad
\tau\to\tau\tau,
\]
and whose columns are the two input labels \(\mathbf 1,\tau\).  The
coefficient choice displayed in \texttt{handoff.md} section 9.6 is the matrix
\[
\eta_{\mathrm{coassoc}}=
\begin{pmatrix}
\varphi^{-1} & 0\\
\varphi^{-1/2} & 0\\
0 & \varphi^{-1}\\
0 & \varphi^{-1}\\
0 & \varphi^{-3/2}
\end{pmatrix}.
\]
It satisfies
\[
\eta_{\mathrm{coassoc}}^{\dagger}\eta_{\mathrm{coassoc}}=I_2.
\]
Lean defines this matrix as \texttt{CoassocBinaryEta} and proves
\texttt{CoassocBinaryEta.conjTranspose * CoassocBinaryEta = 1} in
\texttt{Formalisation/Fibonacci/Coassoc.lean}, using the binary coordinate
isometry theorem in \texttt{Formalisation/Fibonacci/Binary.lean} and the
scalar norm equalities above. Julia checks the four entries of the Gram
matrix numerically at high precision, and WolframScript checks the exact
symbolic identity. AF node \texttt{1.3.17} has been accepted by verifier
\texttt{verifier-coassoc-binary-iso-001}. This is only a fixed-basis
coordinate isometry; it does not prove categorical coassociativity,
associator semantics, fusion-tree basis comparison, or existence of a
categorical map.
\end{theorem}

\section{Renormalization Input and Scaling Operators}

\begin{definition}[Conformal weight input]
The local paper \texttt{references/IsingLikeFibonacciAnyonsKZ.pdf} identifies
the Fibonacci anyon field \(\tau\) with the spin-1 field in the level-3
\(SU(2)\) Wess--Zumino--Witten model, and its table of primary-field
conformal weights gives the spin-1 field the chiral conformal weight
\[
h_\tau=\frac{2}{5}.
\]
For a diagonal theory in which the left and right chiral weights are both
\(h_\tau\), the scaling dimension is their sum:
\[
\Delta_\tau=h_\tau+\overline h_\tau=\frac{2}{5}+\frac{2}{5}=\frac{4}{5}.
\]
This addition convention is sourced locally from
\texttt{references/text/PoilblancOneDimensionalItinerantInteractingAnyons.txt}.
The chiral exponent arithmetic used later for the Fibonacci renormalization
coefficients is
\[
h_{\mathbf 1}-2h_\tau=0-2\cdot\frac25=-\frac45,
\qquad
h_\tau-2h_\tau=\frac25-2\cdot\frac25=-\frac25.
\]
Squaring these powers in the normalisation denominators gives
\[
2\left(h_{\mathbf 1}-2h_\tau\right)=-\frac85,
\qquad
2\left(h_\tau-2h_\tau\right)=-\frac45.
\]
The exact rational equalities are proved in
\texttt{Formalisation/Fibonacci/CFTWeights.lean} and checked in Julia and
WolframScript. AF node \texttt{1.5.5} has been accepted by verifier
\texttt{verifier-cft-weight-001}. AF node \texttt{1.5.7} has been accepted by
verifier \texttt{verifier-cft-rg-exp-001} for the exponent arithmetic. These
formal results do not prove the Wess--Zumino--Witten construction, extraction
of the source table, Wilsonian renormalization coefficients, physical operator
product constants, or a physical renormalization eigenvalue.

The handoff scaling-operator section also uses the rational exponent
convention
\[
\lambda_\tau=b^{-2/5},\qquad
\lambda_{\tau,n}=b^{-(2/5+n)},\qquad
\lambda_{\mathbf 1,n}=b^{-n},\qquad
\lambda_\tau^{\mathrm{diag}}=b^{-4/5}.
\]
Lean proves the corresponding rational exponents
\[
-\frac25,\qquad -\left(\frac25+n\right),\qquad -n,\qquad -\frac45
\]
in \texttt{Formalisation/Fibonacci/CFTWeights.lean}, and Julia and
WolframScript check the same equalities through \(n=6\). AF node
\texttt{1.5.15} has been accepted by verifier
\texttt{verifier-cft-desc-exp-001}. This is only arithmetic for supplied
weights and descendant labels; it does not construct descendant states,
define real powers \(b^{-w}\), prove diagonal renormalization, or construct
physical scaling eigenoperators.
\end{definition}

\begin{definition}[Blocking coefficients]
For a scale factor \(b>1\), set \(\rho=b^{-1}\). The local blocking
coefficients are treated as input data:
\[
\beta^{\mathbf 1}_{\mathbf 1\mathbf 1}=u_0,\quad
\beta^{\mathbf 1}_{\tau\tau}=u_I\rho^{-4/5},
\]
\[
\beta^\tau_{\tau\mathbf 1}=j_L,\quad
\beta^\tau_{\mathbf 1\tau}=j_R,\quad
\beta^\tau_{\tau\tau}=u_\tau\rho^{-2/5}.
\]
The constants \(u_0,u_I,j_L,j_R,u_\tau\) are Wilsonian
renormalization input data; they are not determined by scaling dimensions
alone.
\end{definition}

\begin{theorem}[Fibonacci renormalization amplitude target]
Assuming the coefficients above, the no-mixing polar formula gives
\[
\eta_{\mathbf 1}^{\mathrm{RG}}(\rho)=
\frac{\overline{u_0}v^{\mathbf 1}_{\mathbf 1\mathbf 1}
+\overline{u_I}\rho^{-4/5}v^{\mathbf 1}_{\tau\tau}}
{\sqrt{|u_0|^2+|u_I|^2\rho^{-8/5}}}
\]
and
\[
\eta_\tau^{\mathrm{RG}}(\rho)=
\frac{\overline{j_L}v^\tau_{\tau\mathbf 1}
+\overline{j_R}v^\tau_{\mathbf 1\tau}
+\overline{u_\tau}\rho^{-2/5}v^\tau_{\tau\tau}}
{\sqrt{|j_L|^2+|j_R|^2+|u_\tau|^2\rho^{-4/5}}}.
\]
Status: computational checks in Julia and WolframScript pass for
normalisation and for the finite one-row polar entry formula. The Lean file
\texttt{Formalisation/Fibonacci/RGNoMixing.lean} formalises the two rows as
finite no-mixing vectors
\[
(u_0,u_I s_{\mathrm{vac}}),\qquad
(j_L,j_R,u_\tau s_\tau),
\]
where \(s_{\mathrm{vac}}\) and \(s_\tau\) stand for the supplied powers
\(\rho^{-4/5}\) and \(\rho^{-2/5}\). It proves the displayed denominator
formulas, the normalized-conjugate entries, and row normalisation whenever
the denominator is nonzero. It also proves that the squared norm of the
\(\tau\tau\) amplitude in the tau row is
\[
\frac{|u_\tau|^2s_\tau^2}
{|j_L|^2+|j_R|^2+|u_\tau|^2s_\tau^2},
\]
which becomes the handoff probability
\[
P(\tau\to\tau\tau)=
\frac{|u_\tau|^2\rho^{-4/5}}
{|j_L|^2+|j_R|^2+|u_\tau|^2\rho^{-4/5}}
\]
after the supplied substitution \(s_\tau^2=\rho^{-4/5}\). Lean proves that
this quantity is nonnegative and at most \(1\) when the denominator is
positive. AF node \texttt{1.3.14} has been accepted by
verifier \texttt{verifier-fib-rg-nomix-rows-001}. The scalar no-mixing
normalisation theorem is accepted in AF node \texttt{1.3.7}; the one-row
polar-section formula is accepted in AF node \texttt{1.4.9}; the probability
formula is accepted in AF node \texttt{1.5.16}. The full
polar-decomposition theorem, Wilsonian derivation of the renormalization
coefficients, physical operator product constants, fractional-power analysis,
plotting, physical probability derivation, and categorical basis semantics
remain separate targets.
\end{theorem}

\begin{theorem}[Conditional diagonal scaling]
Let \(I\) be a finite set of scaling labels. For each \(i\in I\), let
\(e_i\) be the coordinate vector with value \(1\) at \(i\) and \(0\)
elsewhere. If a linear map \(R\) is diagonal in this basis with entries
\(\lambda_i\), so that
\[
(Rx)_j=\lambda_j x_j,
\]
then
\[
R e_i=\lambda_i e_i.
\]
In particular, if the entries are stipulated to be
\(\lambda_i=b^{-w_i}\), then \(e_i\) is an eigenvector with eigenvalue
\(b^{-w_i}\). The Lean file
\texttt{Formalisation/LinearAlgebra/DiagonalScaling.lean} proves this
statement, with Julia and WolframScript checks for the chiral pattern
\(w_i=h_i\) and the nonchiral pattern \(w_i=\Delta_i\). Status: AF node
\texttt{1.5.3} is accepted only for this conditional diagonal
linear-algebra statement. The rational scaling exponents for the Fibonacci
primary, descendant, and diagonal nonchiral examples are accepted separately
in AF node \texttt{1.5.15}. These results do not prove real-power analysis,
the diagonalisation of a physical renormalization group map, or the
charge-only negative example.
\end{theorem}

\begin{warning}[Charge-only data do not determine the full scaling spectrum]
The two-dimensional charge-only operator space \(\End(Q)\) does not generally
recover all conformal field theory scaling operators. To recover the full
spectrum, labels must include data such as primary field, descendant level,
and mixing sector.

A concrete finite counterexample is as follows. In the span of two charge
projectors, write
\[
u=(1,1),\qquad c=(1,-1).
\]
For any complex number \(a\), define the charge-only linear map that is
diagonal in the basis \(\{u,c\}\) by
\[
R_a u=u,\qquad R_a c=a\,c.
\]
In the charge-projector coordinates this map is
\[
R_a=
\begin{pmatrix}
(1+a)/2 & (1-a)/2\\
(1-a)/2 & (1+a)/2
\end{pmatrix}.
\]
Thus the charge contrast eigenvalue is the chosen amplitude \(a\). Taking
\(a=1/2\) and \(b=32\) gives \(b^{-2/5}=1/4\), so the charge contrast
eigenvalue is not forced to equal \(b^{-2/5}\). The Lean file
\texttt{Formalisation/LinearAlgebra/ChargeOnly.lean} proves this finite
example, and Julia and WolframScript check it numerically and exactly. Status:
AF node \texttt{1.5.4} is accepted only for this concrete finite
counterexample.
\end{warning}

\section{Temperley--Lieb and Ising Secondary Target}

\begin{definition}[Ising-type fusion rules]
The secondary test category has simple labels
\(\mathbf 1,\sigma,\epsilon\). The local physics source writes the third
label as \(\psi\); in this document \(\epsilon\) means that same label. The
fusion rules are
\[
\sigma\otimes\sigma=\mathbf 1\oplus\epsilon,\quad
\sigma\otimes\epsilon=\sigma,\quad
\epsilon\otimes\epsilon=\mathbf 1.
\]
Together with the unit rule
\[
\mathbf 1\otimes a=a=a\otimes \mathbf 1,
\]
these rules are multiplicity-free except that \(\sigma\otimes\sigma\) has two
simple summands. The local sources are
\texttt{references/text/PenneysUnitaryFusionCategories.md:647},
\texttt{:649}, and
\texttt{references/text/PoilblancOneDimensionalItinerantInteractingAnyons.txt:177}
through \texttt{:188}.
\end{definition}

\begin{theorem}[Finite Ising toy exponent and normalisation]
Use the nonchiral Ising dimensions
\[
\Delta_\sigma=\frac18,\qquad \Delta_\epsilon=1.
\]
Then the exponent for the channel
\[
\sigma\otimes\sigma\to\epsilon
\]
is
\[
\Delta_\epsilon-2\Delta_\sigma
=1-2\cdot\frac18=\frac34.
\]
For a real parameter \(t\), the squared numerator norm of the toy vector with
coefficients
\[
\frac{1}{\sqrt 2},\qquad
\frac{1}{\sqrt 2},\qquad
\frac{t}{2}
\]
is
\[
\left(\frac1{\sqrt2}\right)^2+
\left(\frac1{\sqrt2}\right)^2+
\left(\frac{t}{2}\right)^2
=1+\frac{t^2}{4}.
\]
Thus division by \(\sqrt{1+t^2/4}\) scalar-normalises the toy vector whenever
the source fusion-tree basis vectors are orthonormal. In the handoff
application \(t=\rho^{3/4}\), giving the denominator
\(\sqrt{1+\rho^{3/2}/4}\). The Lean proof is in
\texttt{Formalisation/Ising/Basic.lean}; Julia and WolframScript check the
finite table and rational arithmetic. AF node \texttt{1.8} has been accepted
by verifier \texttt{verifier-ising-toy-001}. This theorem does not construct
the full Ising category, a Temperley--Lieb quotient, Ising \(F\)-symbols,
operator product expansions, a conformal-field-theory spectrum, or a physical
renormalization map.
\end{theorem}

\section{Continuum Structure and Conjectures}

\begin{definition}[Direct-limit continuum Hilbert space]
A coherent family of refinements satisfies
\[
E_{N\to N}=\operatorname{id},\qquad
E_{M\to L}E_{N\to M}=E_{N\to L}.
\]
The continuum charge-sector Hilbert space is the completed direct limit
\[
\mathcal H_\infty^W=\overline{\varinjlim_N\mathcal H_N^W}.
\]
Status: the completed direct-limit Hilbert space remains pending. The finite
coherence algebra for ascending observables and logical lifts is formalised
below.
\end{definition}

\begin{theorem}[Coherent finite ascending channels and logical lifts]
Let \(I\) be a finite coordinate set. Suppose that, for \(n\leq m\), matrices
\[
E_{n,m}\in M_I(\mathbb C)
\]
satisfy
\[
E_{n,n}=I,\qquad E_{m,l}E_{n,m}=E_{n,l}\quad(n\leq m\leq l),
\]
and
\[
E_{n,m}^\dagger E_{n,m}=I.
\]
Define the finite ascending observable map by
\[
A_{n,m}(O)=E_{n,m}^\dagger O E_{n,m}.
\]
Define the finite logical-lift map by
\[
L_{n,m}(O)=E_{n,m} O E_{n,m}^\dagger.
\]
Then
\[
A_{n,m}(I)=I
\]
and, for \(n\leq m\leq l\),
\[
A_{n,m}(A_{m,l}(O))=A_{n,l}(O).
\]
The logical-lift maps satisfy
\[
L_{n,n}(O)=O
\]
and, for \(n\leq m\leq l\),
\[
L_{m,l}(L_{n,m}(O))=L_{n,l}(O).
\]
The Lean file \texttt{Formalisation/LinearAlgebra/CoherentSystem.lean}
proves this finite matrix theorem. Julia and WolframScript check exact
rational rotation examples. AF node \texttt{1.5.6} has been accepted by
verifier \texttt{verifier-coherent-asc-001}, and AF node \texttt{1.5.10} has
been accepted by verifier \texttt{verifier-coherent-lift-001}. This theorem
does not construct a completed direct limit, a continuum Hilbert space, a
continuum algebra, a categorical fine-graining system, or nonconstant growing
Hilbert spaces.
\end{theorem}

\begin{conjecture}[Coherent subdivision family]
A physically correct continuum fine-graining scheme should use a coherent
family
\[
\eta^{(r)}:Q\to Q^{\otimes r},\qquad r=1,2,3,\ldots,
\]
compatible with arbitrary ordered subdivisions.
\end{conjecture}

\begin{conjecture}[Topological fixed points]
For topological fixed points, the coherent family can be expressed using
\(F\)-symbols, quantum dimensions, and local string-net identities.
\end{conjecture}

\begin{conjecture}[Renormalization away from fixed points]
Away from topological fixed points, the coherent family is determined by the
full Wilsonian renormalization trajectory, not by the fusion category alone.
\end{conjecture}

\begin{conjecture}[Universal isometry]
Every physically admissible finite-lattice fine-graining map is, at the
categorical level, an isometry
\[
E_{N\to M}:Q^{\otimes N}\hookrightarrow Q^{\otimes M}.
\]
\end{conjecture}

\begin{conjecture}[Topology and weights]
The fusion category determines which fine-graining channels may be nonzero.
Renormalization data determine the finite-scale coefficients. Isometry then
determines the normalised amplitudes.
\end{conjecture}

\begin{conjecture}[Braiding and order changes]
If a fine-graining interpolation allows particles to exchange spatial order,
then the category must be braided, and the order-change components must use
the braiding. Without braiding, fine-graining maps should preserve order or
avoid crossings.
\end{conjecture}

\section{Current Machine-Checked Artefacts}

\begin{itemize}
\item Lean project: \texttt{lakefile.lean}, \texttt{lean-toolchain}, and
\texttt{Formalisation/}.
\item Lean build command: \texttt{lake build}.
\item Julia check: \texttt{scripts/julia/fibonacci\_checks.jl}.
\item WolframScript check: \texttt{scripts/wolfram/fibonacci\_exact.wls}.
\item Command log: \texttt{results/CHECKS.md}.
\item Master provenance ledger: \texttt{MASTER\_PROVENANCE.md}.
\end{itemize}

\end{document}
